\chapter{Αρχιτεκτονική του Συστήματος}
Το παρόν κεφάλαιο παρουσιάζει τη μεθοδολογία και την αρχιτεκτονική της προτεινόμενης πλατφόρμας Ψηφιακού Διδύμου, με έμφαση στην εξελικτική μετάβασή της από ένα αρχικό τοπικό πρωτότυπο σε ένα ώριμο, \textlatin{cloud-based} οικοσύστημα τύπου \textlatin{Software-as-a-Service (SaaS)}. Η ανάλυση δεν περιορίζεται στην απαρίθμηση των επιμέρους τεχνολογικών συστατικών, αλλά επιχειρεί να τεκμηριώσει τον τρόπο με τον οποίο οι επιλογές υλικού, πρωτοκόλλων, ενδιάμεσου λογισμικού και υπηρεσιών νέφους συνδυάστηκαν σε μία ενιαία, επεκτάσιμη και λειτουργικά συνεκτική αρχιτεκτονική.

Η σχεδιαστική προσέγγιση που ακολουθήθηκε στηρίζεται σε τρεις αλληλένδετες μηχανικές αρχές. Πρώτον, υιοθετείται η λογική της αποσύνδεσης επιπέδων, ώστε η φυσική υποδομή πεδίου, το επίπεδο διαμεσολάβησης, η αποθήκευση και η διεπαφή χρήστη να μπορούν να εξελίσσονται μερικώς ανεξάρτητα. Δεύτερον, διατηρείται η επεξεργασία στην παρυφή, ώστε κρίσιμες μετρήσεις και αυτοματισμοί να μπορούν να λειτουργούν με χαμηλή καθυστέρηση και περιορισμένο δικτυακό φόρτο. Τρίτον, η συνολική πλατφόρμα οργανώνεται ως πραγματικό \textlatin{Digital Twin} και όχι ως απλό \textlatin{telemetry dashboard}, ακολουθώντας τη διάκριση μεταξύ \textlatin{Digital Shadow} και \textlatin{Digital Twin} που προτείνεται στη σχετική βιβλιογραφία \cite{Grieves2017}, \cite{Kritzinger2018}. Η επιλογή αυτή ευθυγραμμίζεται με τη λογική των πολυεπίπεδων \textlatin{Smart Campus} αρχιτεκτονικών \cite{Han2022}.

\section{Χρονοδιάγραμμα και Φάσεις Υλοποίησης}
\subsection{Χρονοδιάγραμμα}
Η ανάπτυξη του προτεινόμενου ψηφιακού διδύμου ακολούθησε μια αυστηρά δομημένη, σταδιακή μεθοδολογία (\textlatin{phased approach}), προκειμένου να διασφαλιστεί η ομαλή μετάβαση από τη θεωρητική μελέτη στην τελική \textlatin{cloud} πλατφόρμα. Το συνολικό χρονοδιάγραμμα του έργου εκτείνεται σε περίοδο 18 μηνών και συνοψίζεται στον Πίνακα~\ref{tab:timeline_ch3}.

\begin{table}[htbp]
\centering
\small
\setlength{\tabcolsep}{4pt}
\renewcommand{\arraystretch}{1.15}
\caption{Χρονοδιάγραμμα ανάπτυξης της πλατφόρμας}
\label{tab:timeline_ch3}
\begin{tabular}{|p{0.22\textwidth}|p{0.24\textwidth}|p{0.42\textwidth}|}
\hline
\textbf{Περίοδος} & \textbf{Φάση} & \textbf{Σύντομη Περιγραφή} \\
\hline
Σεπτέμβριος 2024 & Έναρξη Έρευνας & Εκτενής βιβλιογραφική ανασκόπηση για Ψηφιακά Δίδυμα και \textlatin{Smart Campuses}, εντοπισμός ερευνητικών κενών και ορισμός του σκοπού της μελέτης. \\
\hline
Νοέμβριος 2024 & Σχεδιασμός Αρχιτεκτονικής & Σχεδιασμός της αρχιτεκτονικής, έρευνα και αξιολόγηση εναλλακτικών υλικών (\textlatin{ESP32, Shelly, Sonoff, Raspberry Pi}), επιλογή \textlatin{tech stack} βάσει κριτηρίων κόστους και διαλειτουργικότητας, πρόβλεψη για πολλαπλούς χρήστες και εγκαταστάσεις, και ορισμός μοντέλου ασφάλειας \textlatin{RBAC}. \\
\hline
Μάρτιος 2025 & Πρώτη Επίσκεψη και Εκτίμηση Χώρου & Πραγματοποιήθηκε η πρώτη επιτόπια εκτίμηση των χώρων, η χαρτογράφηση των αναγκών εγκατάστασης και η προετοιμασία της υλοποίησης πεδίου. \\
\hline
Μάιος 2025 & Εγκατάσταση \textlatin{IoT} (\textlatin{Edge Deployment}) & Φυσική εγκατάσταση αισθητήρων, μετρητών ενέργειας, μικροελεγκτών και \textlatin{openHAB} στα Εκπαιδευτήρια Πλάτων και έναρξη συλλογής τηλεμετρίας στις 9 Μαΐου 2025. \\
\hline
Ιούνιος 2025 & Ανάπτυξη Πλατφόρμας (\textlatin{Cloud/SaaS}) & Ανάπτυξη \textlatin{frontend} με \textlatin{React} και \textlatin{Three.js}, ρύθμιση \textlatin{backend} μέσω \textlatin{Supabase} και δημιουργία \textlatin{real-time dashboards}. \\
\hline
Δεκέμβριος 2025 & Ανάλυση Δεδομένων και Τεχνοοικονομική Αξιολόγηση & Επεξεργασία των δεδομένων λειτουργίας (Μάιος--Νοέμβριος 2025), σύγκριση ενεργειακών αποτελεσμάτων με την περίοδο αναφοράς και διαμόρφωση της τεχνοοικονομικής αποτίμησης. \\
\hline
Μάρτιος 2026 & Συγγραφή και Παρουσίαση & Ολοκλήρωση της τεχνικής τεκμηρίωσης, οριστική συγγραφή της διπλωματικής εργασίας και προετοιμασία της τελικής ακαδημαϊκής παρουσίασης. \\
\hline
\end{tabular}
\end{table}

\subsection{Φάσεις Υλοποίησης}
Η τελική μορφή της πλατφόρμας προέκυψε μέσα από μια σταδιακή διαδικασία τεχνολογικής ωρίμανσης. Για τον λόγο αυτό, η αρχιτεκτονική της εργασίας παρουσιάζεται ως ακολουθία διακριτών φάσεων υλοποίησης και όχι ως στιγμιαία, μονοσήμαντη σχεδίαση.

\noindent\textbf{Φάση Α: Ανάλυση Απαιτήσεων και Χαρτογράφηση Συστήματος.} Κατά την πρώτη φάση πραγματοποιήθηκε η αποτύπωση των λειτουργικών απαιτήσεων του σχολικού συγκροτήματος, η επιλογή των κρίσιμων χώρων προς παρέμβαση και η αξιολόγηση εναλλακτικών υλικών βάσει κριτηρίων κόστους, διαθεσιμότητας και διαλειτουργικότητας με ανοιχτό λογισμικό, και η χαρτογράφηση των βασικών ροών δεδομένων. Σε αυτό το στάδιο ορίστηκαν οι κρίσιμες μεταβλητές της μελέτης, όπως ο θόρυβος στις αίθουσες, η ενεργειακή κατανάλωση βαρέων φορτίων, η θερμοκρασία εντός των ψυκτικών θαλάμων και η κατάσταση υποδομών πρόσβασης. Παράλληλα, ορίστηκαν οι βασικοί επιχειρησιακοί ρόλοι του συστήματος: παρακολούθηση, ειδοποίηση, απομακρυσμένος έλεγχος και συγκεντρωτική εποπτεία.

\noindent\textbf{Φάση Β: Τοπικό Πρωτότυπο Υλικού και \textlatin{openHAB}.} Η δεύτερη φάση αντιστοιχεί στο αρχικό \textlatin{proof of concept} της πλατφόρμας. Σε αυτήν εγκαταστάθηκε το \textlatin{edge hardware}, ενεργοποιήθηκε η πρώτη αξιόπιστη ροή τηλεμετρίας μέσω \textlatin{MQTT} και υλοποιήθηκαν οι πρώτοι κανόνες \textlatin{closed-loop} ελέγχου στο \textlatin{openHAB}. Το στάδιο αυτό επέτρεψε τη δοκιμή της λογικής των αυτοματισμών, τη βαθμονόμηση των αισθητήρων, την καταγραφή ιστορικών δεδομένων και την πρώτη πρακτική επιβεβαίωση ότι το προτεινόμενο σύστημα μπορεί να λειτουργήσει ως \textlatin{Digital Twin} και όχι μόνο ως \textlatin{Digital Shadow}.

\noindent\textbf{Φάση Γ: Ανάπτυξη \textlatin{Cloud SaaS} Διεπαφής.} Η τρίτη φάση αφορά τη μετάβαση από το πρωτότυπο τοπικό περιβάλλον ελέγχου σε αρχιτεκτονική \textlatin{cloud SaaS}. Η μετάβαση αυτή ήταν αναγκαία, επειδή οι αρχικές διεπαφές του πρωτοτύπου εξυπηρετούσαν κυρίως τεχνικούς χρήστες και τοπικά σενάρια δοκιμών, ενώ η νέα αρχιτεκτονική διαχωρίζει το επίπεδο συλλογής και αυτοματισμού από το επίπεδο διάθεσης υπηρεσιών και διεπαφής χρήστη. Με τον τρόπο αυτό η πλατφόρμα καθίσταται προσβάσιμη και σε μη τεχνικούς ρόλους, όπως διευθυντές και διοικητικό προσωπικό, και μπορεί να επεκταθεί σε περισσότερους χρήστες και εγκαταστάσεις χωρίς μεταβολή του φυσικού επιπέδου.

\noindent\textbf{Φάση Δ: Ολοκλήρωση, Δοκιμές και Επιχειρησιακή Ωρίμανση.} Στην τελική φάση πραγματοποιήθηκε η ολοκλήρωση του \textlatin{edge middleware} με το \textlatin{cloud SaaS} επίπεδο, η επαλήθευση των ροών δεδομένων και η σταδιακή ενσωμάτωση πρόσθετων λειτουργιών, όπως πίνακες ελέγχου και \textlatin{2D/3D} οπτικοποιήσεις. Το αποτέλεσμα αυτής της φάσης είναι μια αρχιτεκτονική στην οποία το τοπικό επίπεδο διατηρεί τις λειτουργίες συλλογής και αυτοματισμού, ενώ το \textlatin{cloud} επίπεδο αναλαμβάνει την αποθήκευση, την παρουσίαση των δεδομένων, τη διαχείριση ρόλων και την απομακρυσμένη πρόσβαση. Ο διαχωρισμός αυτός διευκολύνει τη συντήρηση και τη μελλοντική επέκταση της πλατφόρμας χωρίς ανασχεδιασμό του υλικού πεδίου.

\section{Μελέτη Περίπτωσης και Ανάλυση Απαιτήσεων}
Η αρχιτεκτονική της πλατφόρμας διαμορφώθηκε πάνω στις πραγματικές επιχειρησιακές ανάγκες των Εκπαιδευτηρίων Πλάτων, τα οποία επιλέχθηκαν ως πεδίο εφαρμογής της έρευνας. Το σχολικό συγκρότημα συγκεντρώνει χαρακτηριστικά που το καθιστούν αντιπροσωπευτικό παράδειγμα σύγχρονης εκπαιδευτικής υποδομής: πολλαπλούς λειτουργικούς χώρους, συνεχόμενη καθημερινή χρήση, διακριτά ενεργειακά φορτία, και ανάγκη για κεντρικό έλεγχο χωρίς να διαταράσσεται η εκπαιδευτική διαδικασία.

Η αρχική ανάλυση απαιτήσεων ανέδειξε τρεις διακριτούς άξονες σχεδιασμού. Ο πρώτος άξονας αφορούσε το μαθησιακό περιβάλλον και ειδικότερα την ακουστική άνεση εντός των αιθουσών. Η απαίτηση δεν ήταν απλώς η καταγραφή της στάθμης θορύβου, αλλά η δυνατότητα ήπιας και μη παρεμβατικής ανατροφοδότησης προς την τάξη όταν η στάθμη υπερέβαινε ένα επιχειρησιακό όριο παρέμβασης. Ο δεύτερος άξονας αφορούσε την ενεργειακή και λειτουργική διαχείριση του σχολείου, με έμφαση στο κυλικείο, όπου εντοπίστηκαν ενεργοβόρα φορτία, ψυκτικός εξοπλισμός και ανάγκη έγκαιρης ειδοποίησης για λειτουργικές ανωμαλίες. Ο τρίτος άξονας αφορούσε τη διαλειτουργικότητα με υφιστάμενες υποδομές, όπως οι αυτόματες πύλες του σχολείου, ώστε η αρχιτεκτονική να μην περιορίζεται σε μία νέα αυτόνομη εγκατάσταση, αλλά να μπορεί να ενοποιεί και \textlatin{legacy} συστήματα μέσα από λογισμικό.

Οι απαιτήσεις αυτές οδήγησαν σε μια αρχιτεκτονική που έπρεπε ταυτόχρονα να είναι χαμηλού κόστους στο φυσικό επίπεδο, επεκτάσιμη στο επίπεδο λογισμικού, κατανοητή για μη τεχνικούς χρήστες και συμβατή με τη λογική ενός σύγχρονου \textlatin{Smart Campus}. Κατά συνέπεια, το σύστημα δεν σχεδιάστηκε ως μία απλή εφαρμογή αυτοματισμού, αλλά ως ένα ενιαίο περιβάλλον παρακολούθησης, ελέγχου και επιχειρησιακής εποπτείας.

\section{Σχεδιασμός Υλικού Εξοπλισμού και Επιλογή Συσκευών}
Ο φυσικός εξοπλισμός της πλατφόρμας διατηρεί τον \textlatin{low-cost} και \textlatin{open} χαρακτήρα του, καθώς αυτός αποτελεί βασική προϋπόθεση για τη δυνατότητα επανάληψης της αρχιτεκτονικής σε σχολικά περιβάλλοντα με περιορισμένους πόρους. Αντί της υιοθέτησης μιας κλειστής, βαριάς και υψηλού κόστους υποδομής τύπου \textlatin{BMS}, επιλέχθηκε ένα σύνολο εμπορικά διαθέσιμων συσκευών με ευρεία κοινότητα υποστήριξης, ώριμα ανοικτά πρωτόκολλα και δυνατότητα διασύνδεσης με \textlatin{middleware} και \textlatin{cloud layers}.

\subsection{Κόμβοι Ακουστικής Παρακολούθησης}
Για την παρακολούθηση του θορύβου στις αίθουσες διδασκαλίας επιλέχθηκαν μικροελεγκτές \textlatin{ESP32}, οι οποίοι συνδυάζουν χαμηλό κόστος, επαρκή υπολογιστική ισχύ και εγγενή υποστήριξη δικτύωσης \textlatin{Wi-Fi}. Η επιλογή τους δικαιολογείται από το γεγονός ότι μπορούν να εκτελούν τοπική επεξεργασία σημάτων και να υπολογίζουν μετρικές όπως η ισοδύναμη συνεχής στάθμη ηχητικής πίεσης (\textlatin{LAeq}), χωρίς να απαιτείται μεταφορά ακατέργαστου σήματος προς το κεντρικό σύστημα. Σε επίπεδο αισθητήρα χρησιμοποιήθηκε ψηφιακό μικρόφωνο \textlatin{I2S}, όπως το \textlatin{INMP441}, το οποίο προσφέρει επαρκή ποιότητα δειγματοληψίας για τη συγκεκριμένη εφαρμογή.

Σημαντική διευκρίνιση ως προς την προστασία προσωπικών δεδομένων: ο αισθητήρας \textlatin{INMP441} δεν πραγματοποιεί ηχογράφηση και δεν αποθηκεύει κανένα ηχητικό περιεχόμενο. Το μόνο που μετράται και αποθηκεύεται είναι ο αριθμητικός δείκτης στάθμης θορύβου (\textlatin{LAeq} σε \textlatin{dBA}) --- δηλαδή η ένταση του ήχου, όχι το περιεχόμενό του. Αυτό καθιστά το σύστημα συμβατό με τον Γενικό Κανονισμό Προστασίας Δεδομένων (\textlatin{GDPR}/\textlatin{ΕΕ 2016/679}), καθώς δεν συλλέγονται προσωπικά δεδομένα ήχου που θα μπορούσαν να χρησιμοποιηθούν για αναγνώριση ή παρακολούθηση προσώπων.

Η επιλογή αυτής της διάταξης δεν υπαγορεύθηκε μόνο από οικονομικά κριτήρια. Εξίσου σημαντικό ήταν το γεγονός ότι οι κόμβοι μπορούν να ενσωματωθούν εύκολα στο συνολικό \textlatin{MQTT-based} οικοσύστημα της πλατφόρμας, να αναβαθμιστούν λογισμικά και να επεκταθούν σε πρόσθετες αίθουσες χωρίς εκ βάθρων αλλαγές στην αρχιτεκτονική.

\subsection{Ενεργειακοί Μετρητές και Αισθητήρες Κυλικείου}
Για τη μέτρηση της ενεργειακής κατανάλωσης επιλέχθηκαν έξυπνοι μετρητές \textlatin{Shelly 1PM}, καθώς παρέχουν δυνατότητα καταγραφής ισχύος, τάσης και ενέργειας σε πραγματικό χρόνο, με παράλληλη ευχέρεια ενσωμάτωσης σε ανοικτά πρωτόκολλα και \textlatin{middleware workflows}. Η επιλογή της συγκεκριμένης οικογένειας συσκευών δικαιολογείται από τη χαμηλή τιμή κτήσης, την ωριμότητα του οικοσυστήματος και τη δυνατότητα εγκατάστασης σε υφιστάμενα κυκλώματα χωρίς ριζική ανακατασκευή της ηλεκτρολογικής υποδομής.

Η μόνη παρακολούθηση θερμοκρασίας που υλοποιήθηκε στο σύστημα αφορά το ψυγείο του κυλικείου, όπου προστέθηκαν αισθητήρες θερμοκρασίας \textlatin{DS18B20} και μαγνητικές επαφές θυρών, ώστε να συνδυαστούν ηλεκτρικές και λειτουργικές μετρήσεις. Η σύζευξη αυτών των δεδομένων επιτρέπει την ανίχνευση όχι μόνο του ενεργειακού προφίλ του εξοπλισμού, αλλά και της πιθανής αιτίας αποκλίσεων, όπως παρατεταμένο άνοιγμα της θύρας ή ασυνήθιστα αυξημένος χρόνος λειτουργίας του συμπιεστή.

\subsection{Ενεργοποιητές και Συστήματα Παρέμβασης}
Στις αίθουσες χρησιμοποιήθηκαν ενεργοποιητές ρελέ \textlatin{Sonoff Basic R2} για τον έλεγχο της \textlatin{quiet-lamp}, ενώ στο επίπεδο του κυλικείου χρησιμοποιήθηκε ηχητικός ενεργοποιητής (\textlatin{buzzer}) συνδεδεμένος με αισθητήρα μαγνητικής επαφής θύρας, ώστε να παράγεται τοπική ειδοποίηση όταν η πόρτα του ψυγείου παραμένει ανοιχτή πέραν του επιτρεπτού χρονικού ορίου (45 δευτερόλεπτα). Οι ενεργοποιητές αυτοί επιλέχθηκαν λόγω χαμηλού κόστους, ευκολίας ενσωμάτωσης και υποστήριξης λογικής \textlatin{software-defined control}. Με τον τρόπο αυτό η πλατφόρμα δεν περιορίζεται σε \textlatin{monitoring}, αλλά επεκτείνεται σε πραγματική παρέμβαση στο φυσικό περιβάλλον. Το Σχήμα~\ref{fig:field_installation} παρουσιάζει στιγμιότυπα από την εγκατάσταση των παραπάνω συσκευών στο πεδίο.

\begin{figure}[htbp]
\centering
\includegraphics[angle=-90,width=0.31\textwidth]{fig_3_6}\hfill
\includegraphics[angle=-90,width=0.31\textwidth]{fig_3_7}\hfill
\includegraphics[angle=-90,width=0.31\textwidth]{fig_3_8}

\vspace{0.5em}

\includegraphics[angle=-90,width=0.31\textwidth]{fig_3_9}\hfill
\includegraphics[angle=-90,width=0.31\textwidth]{fig_3_10}\hfill
\includegraphics[angle=-90,width=0.31\textwidth]{fig_3_11}

\vspace{0.5em}

\caption{Φωτογραφίες από την εγκατάσταση στα Εκπαιδευτήρια Πλάτων. Πάνω σειρά: εγκατάσταση κόμβων ακουστικής παρακολούθησης και \textlatin{quiet-lamp} στις αίθουσες διδασκαλίας. Κάτω σειρά: εγκατάσταση αισθητήρων και έξυπνων μετρητών στον χώρο του κυλικείου.}
\label{fig:field_installation}
\end{figure}

\section{Η Αρχιτεκτονική των Πέντε Επιπέδων}
Η αρχιτεκτονική της πλατφόρμας υιοθετεί και προσαρμόζει το θεωρητικό μοντέλο πέντε επιπέδων της βιβλιογραφίας στις απαιτήσεις του σχολικού περιβάλλοντος, όπως αποτυπώνεται στο Σχήμα~\ref{fig:five_layer_architecture}. Τα επίπεδα \textlatin{Physical} και \textlatin{Connection} παραμένουν ως έχουν, ενώ τα επίπεδα \textlatin{Data}, \textlatin{Virtual} και \textlatin{Service} προσαρμόζονται σε \textlatin{Edge Middleware}, \textlatin{Cloud Data and Services} και \textlatin{Application and UX} αντίστοιχα --- αντικατοπτρίζοντας την πραγματική τεχνολογική στοίβα της υλοποίησης.

\begin{figure}[htbp]
\centering
\includegraphics[width=\textwidth]{fig_3_1}
\caption{Συνοπτική απεικόνιση της αρχιτεκτονικής πέντε επιπέδων της πλατφόρμας Ψηφιακού Διδύμου.}
\label{fig:five_layer_architecture}
\end{figure}

\subsection{\texorpdfstring{\textlatin{Physical Layer}}{Physical Layer}}
Το \textlatin{Physical Layer} περιλαμβάνει το σύνολο του υλικού εξοπλισμού της εγκατάστασης: τους κόμβους \textlatin{ESP32/INMP441} για μέτρηση ακουστικής στάθμης στις αίθουσες, τον μετρητή ισχύος \textlatin{Shelly 1PM Mini Gen3} για παρακολούθηση κατανάλωσης ψυγείου, τον κόμβο \textlatin{ESP8266} με μαγνητική επαφή \textlatin{MC-38} για ανίχνευση ανοίγματος θύρας, τα ρελέ \textlatin{Sonoff Basic R2} για έλεγχο \textlatin{quiet-lamp} και \textlatin{buzzer}, και τις \textlatin{dry-contact} επαφές ενοποίησης με τον υπάρχοντα ελεγκτή των αυτόματων πυλών. Όλες οι συσκευές επικοινωνούν ασύρματα μέσω \textlatin{WiFi} με το επόμενο επίπεδο.

\subsection{\texorpdfstring{\textlatin{Connection Layer}}{Connection Layer}}
Το \textlatin{Connection Layer} αναλαμβάνει τη μεταφορά δεδομένων μεταξύ φυσικών συσκευών και λογισμικού. Ο \textlatin{MQTT broker} \textlatin{Mosquitto} εκτελείται στον \textlatin{Raspberry Pi 4} και δέχεται \textlatin{telemetry} από όλους τους κόμβους. Τα \textlatin{openHAB Bindings} (\textlatin{MQTT Binding}, \textlatin{HTTP Binding}) μεταφράζουν τα μηνύματα σε \textlatin{openHAB Items}, δημιουργώντας ενιαία αναπαράσταση ανεξάρτητα από το πρωτόκολλο κάθε συσκευής.

\subsection{\texorpdfstring{\textlatin{Edge Middleware Layer}}{Edge Middleware Layer}}
Το \textlatin{Edge Middleware Layer} αποτελεί τον εγκέφαλο της τοπικής υποδομής και εκτελείται εξ ολοκλήρου στον \textlatin{Raspberry Pi 4}. Το \textlatin{openHAB 4.x Core} με το \textlatin{Semantic Model} (\textlatin{Location}$\rightarrow$\textlatin{Building}$\rightarrow$\textlatin{Floor}$\rightarrow$\textlatin{Room}$\rightarrow$\textlatin{Equipment}$\rightarrow$\textlatin{Point} ιεραρχία) οργανώνει σημασιολογικά όλες τις συσκευές της εγκατάστασης. Ο \textlatin{Rules Engine} υλοποιεί τη λογική αυτοματισμού --- τον κανόνα \textlatin{NoiseGuard} για ακουστική παρέμβαση και τον κανόνα \textlatin{FridgeGuardian} για παρακολούθηση ψυγείου --- κλείνοντας τον βρόχο ελέγχου από το \textlatin{Sense} στο \textlatin{Act}. Η \textlatin{InfluxDB} αποθηκεύει χρονοσειρές τηλεμετρίας για ιστορική ανάλυση.

\subsection{\texorpdfstring{\textlatin{Cloud Data and Services Layer}}{Cloud Data and Services Layer}}
Το \textlatin{Cloud Data and Services Layer} παρέχει την απομακρυσμένη αποθήκευση και τις υπηρεσίες \textlatin{backend} της πλατφόρμας. Η \textlatin{Supabase PostgreSQL} αποθηκεύει τη διαμόρφωση της πλατφόρμας, τα δεδομένα χρηστών και το ιστορικό διαμόρφωσης και \textlatin{installation metadata}. Οι \textlatin{Supabase Edge Functions} παρέχουν \textlatin{REST API} για επικοινωνία με το \textlatin{frontend} και ενσωμάτωση με το \textlatin{Gemini API} για \textlatin{LLM-based energy forecasting}. Το σύστημα \textlatin{real-time alerts} αποστέλλει αυτόματες ειδοποιήσεις σε ανωμαλίες που εντοπίζει ο \textlatin{Rules Engine}.

\subsection{\texorpdfstring{\textlatin{Application and UX Layer}}{Application and UX Layer}}
Το \textlatin{Application and UX Layer} παρέχει την τελική διεπαφή προς τον χρήστη μέσω της \textlatin{Cloud SaaS} πλατφόρμας που είναι προσβάσιμη στη διεύθυνση \textlatin{digitaltwin.pi-tech.gr}. Το \textlatin{React 18 dashboard} προσφέρει \textlatin{3D campus visualization} μέσω \textlatin{React Three Fiber}, \textlatin{2D SVG floor plans} με κατάσταση αιθουσών, \textlatin{real-time} παρακολούθηση αισθητήρων και αμφίδρομο έλεγχο συσκευών. Η πρόσβαση διέπεται από \textlatin{RBAC} πολυεπίπεδης εξουσιοδότησης με ξεχωριστούς ρόλους για διαχειριστή και χρήστη.

\section{Σημασιολογική Οργάνωση και Διαλειτουργικότητα}
Η παρούσα πλατφόρμα σχεδιάστηκε με έμφαση στη διαλειτουργικότητα, ώστε να μην περιορίζεται σε νέες συσκευές \textlatin{IoT}, αλλά να μπορεί να ενοποιεί και υφιστάμενες εγκαταστάσεις του σχολείου. Χαρακτηριστικό παράδειγμα αποτελεί η διασύνδεση με το ήδη εγκατεστημένο σύστημα αυτόματων πυλών, το οποίο εντάχθηκε στο ενιαίο περιβάλλον ελέγχου χωρίς αντικατάσταση του φυσικού εξοπλισμού, αλλά μέσω λογισμικής διασύνδεσης και κοινής διαχείρισης κατάστασης.

Η επιλογή αυτή ακολουθεί τη λογική του \textlatin{Brownfield Integration}: αντί για εκ βάθρων ανακατασκευή της υποδομής, η πλατφόρμα λειτουργεί ως επίπεδο ενοποίησης ανάμεσα σε ετερογενή υποσυστήματα, επιτρέποντας κοινή εποπτεία, ενοποιημένους κανόνες ελέγχου και σταδιακή επέκταση σε νέα σημεία του σχολείου. Από τεχνική άποψη, η προσέγγιση αυτή μειώνει το κόστος μετάβασης, περιορίζει την εξάρτηση από ιδιόκτητα συστήματα και καθιστά ρεαλιστική την υιοθέτηση του \textlatin{Digital Twin} σε υφιστάμενα σχολικά κτίρια.




